\uuid{Q3pH}
\titre{Distance}
\auteur{bodin ; hanani ; mzari}
\organisation{exo7}
\licence{CC BY-NC-SA 4.0}
\datecreate{2020-01-20}
\niveau{L1}
\module{Géométrie}
\chapitre{Géométrie de l'espace}
\sousChapitre{Distance}
\difficulte{2}
\theme{Distance}
\exo7qcmid{lille-1.geometrie-espace.024}

\contenu{
\texte{
Pour ces questions, l'espace est muni d'un repère orthonormé direct  $(O, \vec{i}, \vec{j}, \vec{k})$.
}

\begin{question}
\qclassification{141.05}
\qtags{motcle=distance}
On considère les droites $D_1 : \left\{\begin{array}{ccl}x&=&1+t\\y&=&-t\\ z&=&1+t,\quad (t \in \Rr)  \end{array}\right.$ et
$D_2 : \left\{\begin{array}{ccl}y&=&2\\x-z&=&2 \end{array}\right.$
La distance entre $D_1$ et $D_2$ est :
\begin{answers}

    \bad {$0$}

    \bad {$\frac{1}{\sqrt 2}$}

   \good{$\sqrt 2$}

    \bad{Les droites se rapprochent autant que l'on veut sans se toucher.}

\end{answers}
\begin{explanations}
Si $D_1$ passe par un point $A_1$ et est dirigée par un vecteur $\vec{u_1}$ et
 $D_2$ passe par un point $A_2$ et est dirigée par un vecteur $\vec{u_2}$, la distance entre $D_1$ et $D_2$ est donnée par : $\frac{|\det(\overrightarrow{A_1A_2},\vec{u_1}, \vec{u_2})|}{\Vert \vec{u_1} \wedge \vec{u_2} \Vert }$.
\end{explanations}
\end{question}
}
